\documentclass[11pt,a4paper]{article}
\usepackage[T1]{fontenc}
\usepackage[utf8]{inputenc}
\usepackage[brazilian]{babel}
\usepackage[default,type1,semibold]{inter}
\usepackage[a4paper,top=18mm,bottom=24mm,left=18mm,right=18mm,footskip=12mm]{geometry}
\usepackage[table]{xcolor}
\usepackage{amssymb,textcomp,array,longtable,enumitem,fancyhdr,lastpage,needspace,fontawesome5}
\usepackage[skins,breakable,raster]{tcolorbox}
\usepackage[hidelinks]{hyperref}
\usetikzlibrary{positioning,calc,babel}
\DeclareUnicodeCharacter{2192}{\ensuremath{\rightarrow}}
\DeclareUnicodeCharacter{2212}{\ensuremath{-}}
\DeclareUnicodeCharacter{2265}{\ensuremath{\geq}}
\DeclareUnicodeCharacter{2248}{\ensuremath{\approx}}
\DeclareUnicodeCharacter{03C3}{\ensuremath{\sigma}}

\definecolor{cor}{HTML}{7C3AED}
\definecolor{suave}{HTML}{F1EBFE}
\definecolor{texto}{HTML}{0F172A}
\definecolor{cinza}{HTML}{5B6B84}
\definecolor{linha}{HTML}{D3DAE6}
\color{texto}
\setlength{\parindent}{0pt}
\setlength{\parskip}{4pt}
\linespread{1.1}
\setlength{\LTpre}{6pt}
\setlength{\LTpost}{10pt}
\renewcommand{\arraystretch}{1.45}
\arrayrulecolor{linha}
\setlist{itemsep=4pt,topsep=4pt,parsep=0pt}
\setlist[itemize]{label={\color{cor}\small\faCaretRight},leftmargin=1.4em}
% Passos numerados: número dentro de um círculo da cor da equipe.
\DeclareRobustCommand{\bolinha}[1]{\tikz[baseline=(n.base)]\node[circle,fill=cor,text=white,inner sep=1.6pt,minimum size=17pt,font=\footnotesize\bfseries](n){#1};}
\setlist[enumerate]{label=\bolinha{\arabic*},leftmargin=2.3em,labelsep=0.6em,itemsep=6pt}
% Checklist: quadrado grande para marcar à mão.
\DeclareRobustCommand{\caixinha}{\tikz[baseline=-0.6ex]\draw[cor,line width=1.2pt,rounded corners=2pt](0,-0.16)rectangle(0.42,0.26);}
\newlist{checklist}{itemize}{1}
\setlist[checklist]{label=\caixinha,leftmargin=2em,labelsep=0.6em,itemsep=7pt}

\pagestyle{fancy}
\fancyhf{}
\renewcommand{\headrulewidth}{0pt}
\fancyfoot[L]{\footnotesize\color{cor}\textbf{\faRunning\ \ Frente B}\color{cinza}\ \ ·\ \ Produto e Tarefas (PO + Scrum Master) · Passo a passo}
\fancyfoot[R]{\footnotesize\color{cinza}\textbf{\thepage}\,/\,\pageref{LastPage}}
% Fita da cor da equipe na margem esquerda de todas as páginas.
\fancyhead[L]{\begin{tikzpicture}[remember picture,overlay]\fill[cor](current page.north west)rectangle([xshift=5pt]current page.south west);\end{tikzpicture}}

% Título de seção com ícone: nunca sozinho no pé da página.
\newcommand{\secao}[2]{\Needspace{10\baselineskip}\vspace{12pt}%
  \noindent\tikz[baseline=(i.base)]\node[circle,fill=cor,text=white,minimum size=26pt,inner sep=0pt](i){\small #1};%
  \hspace{8pt}{\Large\bfseries #2}\par\vspace{-4pt}{\color{cor!35}\rule{\linewidth}{1.2pt}}\par\vspace{2pt}}
\newcommand{\subsecao}[1]{\Needspace{8\baselineskip}\vspace{8pt}{\large\bfseries\color{cor}#1}\par}
\newcommand{\rotulo}[2]{\Needspace{6\baselineskip}\vspace{10pt}%
  \noindent{\color{cor}\normalsize #1}\hspace{6pt}{\footnotesize\bfseries\color{cor}\MakeUppercase{#2}}\par\vspace{2pt}}
\newtcolorbox{destaque}[1][]{enhanced,breakable,colback=suave,colframe=suave,borderline west={4pt}{0pt}{cor},arc=4pt,boxrule=0pt,left=12pt,right=12pt,top=8pt,bottom=8pt,before skip=8pt,after skip=10pt,#1}
\newtcolorbox{apostila}{enhanced,breakable,colback=white,colframe=cor,boxrule=1pt,arc=6pt,left=12pt,right=12pt,top=8pt,bottom=8pt,before skip=8pt,after skip=8pt,
  overlay={\node[fill=cor,text=white,rounded corners=3pt,font=\scriptsize\bfseries,inner sep=4pt,anchor=west] at ([xshift=10pt]frame.north west) {\faBookOpen\ DA APOSTILA};}}
\newtcolorbox{selo}{enhanced,left=8pt,right=8pt,top=8pt,bottom=8pt,before skip=0pt,after skip=0pt,
  frame code={\draw[cor,dashed,line width=1.4pt,rounded corners=8pt](frame.south west)rectangle(frame.north east);},
  interior code={\fill[suave,rounded corners=8pt](interior.south west)rectangle(interior.north east);}}
% Anotações: ocupa o resto da página da etapa (só se sobrar espaço útil).
\newcommand{\anotacoes}{\par\ifdim\dimexpr\pagegoal-\pagetotal\relax>3.5cm\ifdim\pagegoal<\maxdimen
\begin{tcolorbox}[enhanced,height fill,colback=white,colframe=cor!35,boxrule=0.8pt,arc=8pt,before skip=10pt,left=10pt,top=6pt,
  borderline={0.8pt}{0pt}{cor!35,dashed},frame hidden,
  underlay={\begin{scope}\clip(interior.south west)rectangle(interior.north east);
    \foreach \y in {1,...,40}{\draw[cor!18,line width=0.6pt]([yshift=-\y*0.75cm-0.5cm]interior.north west)--([yshift=-\y*0.75cm-0.5cm]interior.north east);}\end{scope}}]
{\footnotesize\bfseries\color{cor}\faStickyNote\ \ ANOTAÇÕES}
\end{tcolorbox}\fi\fi}
% Fim da etapa: se o bloco final não cabe, a sobra vira Anotações e o bloco vai para a próxima página.
\newif\ifjaanotou
\newsavebox{\blocofinal}
% Mede o bloco final já montado; se não couber no resto da página, a sobra vira Anotações.
\newcommand{\usarblocofinal}{\par\global\jaanotoufalse
  \ifdim\pagegoal<\maxdimen\ifdim\dimexpr\pagegoal-\pagetotal\relax<\dimexpr\ht\blocofinal+\dp\blocofinal+18pt\relax
    \anotacoes\global\jaanotoutrue\clearpage\fi\fi
  \vspace{10pt}\noindent\usebox{\blocofinal}\par}
\newtcolorbox{faixa}[1][]{enhanced,colback=cor,colframe=cor,coltext=white,arc=8pt,boxrule=0pt,left=14pt,right=14pt,top=10pt,bottom=10pt,#1}
\newtcbox{\chip}[1][]{enhanced,on line,colback=suave,colframe=suave,coltext=cor,arc=7pt,boxrule=0pt,left=5pt,right=5pt,top=2pt,bottom=2pt,fontupper=\footnotesize\bfseries,#1}
\newtcbox{\chipbranco}{enhanced,on line,colback=white,colframe=white,coltext=cor,arc=7pt,boxrule=0pt,left=6pt,right=6pt,top=2pt,bottom=2pt,fontupper=\footnotesize\bfseries}
\begin{document}
\thispagestyle{fancy}
\begin{tikzpicture}[remember picture,overlay]
\fill[cor,opacity=0.10] ([xshift=-1.5cm,yshift=-1cm]current page.north east) circle (5.2cm);
\fill[cor,opacity=0.07] ([xshift=-5.5cm,yshift=-4.5cm]current page.north east) circle (2.4cm);
\fill[cor,opacity=0.06] ([xshift=2cm,yshift=3cm]current page.south west) circle (4cm);
\end{tikzpicture}
\vspace*{2mm}
\noindent\begin{tikzpicture}
\node[circle,fill=cor,text=white,minimum size=78pt,font=\fontsize{40}{40}\selectfont\bfseries] (L) {B};
\node[right=14pt of L.east,anchor=west,align=left,text width=12.2cm] {
  {\footnotesize\bfseries\color{cor}PROJETO ESTADIA\ \ ·\ \ EQUIPE VIOLETA\ \ ·\ \ FRENTE B}\\[4pt]
  {\fontsize{23}{27}\selectfont\bfseries Produto e Tarefas (PO + Scrum Master)}\\[6pt]
  \chip{\faShoePrints\ \ Passo a passo}\ \chip{\faRunning\ \ Aulas 01, 02, 04 e 07 da apostila}
};
\end{tikzpicture}
\vspace{10pt}

\begin{faixa}
{\large\faRocket}\hspace{8pt}{\large\bfseries Nossa missão}\par\vspace{2pt}
{\normalsize Organizar a lista do que o sistema precisa ter e conduzir as tarefas do grupo do jeito ágil (Scrum e Kanban).}
\end{faixa}

\vspace{6pt}
\begin{tcolorbox}[enhanced,colback=white,colframe=cor!40,boxrule=1pt,arc=8pt,left=10pt,right=10pt,top=8pt,bottom=10pt,
  title={\faUserFriends\ \ Nossa equipe},fonttitle=\bfseries\small,coltitle=cor,colbacktitle=suave,attach boxed title to top left={xshift=10pt,yshift=-8pt},boxed title style={arc=4pt,boxrule=0pt}]
\vspace{6pt}
\begin{minipage}[t]{0.31\linewidth}\centering
\tikz\node[circle,draw=cor,line width=1.4pt,fill=white,minimum size=34pt,text=cor]{\large \faCrown};\par\vspace{4pt}
{\scriptsize\bfseries\color{cor}\MakeUppercase{Product Owner}}\par\vspace{14pt}
{\color{cor!60}\rule{0.9\linewidth}{0.8pt}}\par{\scriptsize\color{cinza}nome}
\end{minipage}\hfill
\begin{minipage}[t]{0.31\linewidth}\centering
\tikz\node[circle,draw=cor,line width=1.4pt,fill=white,minimum size=34pt,text=cor]{\large \faRunning};\par\vspace{4pt}
{\scriptsize\bfseries\color{cor}\MakeUppercase{Scrum Master}}\par\vspace{14pt}
{\color{cor!60}\rule{0.9\linewidth}{0.8pt}}\par{\scriptsize\color{cinza}nome}
\end{minipage}\hfill
\begin{minipage}[t]{0.31\linewidth}\centering
\tikz\node[circle,draw=cor,line width=1.4pt,fill=white,minimum size=34pt,text=cor]{\large \faLaptopCode};\par\vspace{4pt}
{\scriptsize\bfseries\color{cor}\MakeUppercase{Dev de apoio}}\par\vspace{14pt}
{\color{cor!60}\rule{0.9\linewidth}{0.8pt}}\par{\scriptsize\color{cinza}nome}
\end{minipage}
\end{tcolorbox}
\secao{\faMapSigns}{Mapa da jornada}
{\color{cinza}4 tarefas curtas, uma por página. Marquem a bolinha quando terminarem cada uma.}\par\vspace{4pt}
\begin{tcolorbox}[enhanced,colback=suave!60,colframe=cor!25,boxrule=0.6pt,arc=6pt,left=6pt,right=8pt,top=4pt,bottom=4pt,before skip=0pt,after skip=4pt]
\begin{tabular}{@{}m{24pt}m{20pt}>{\raggedright\arraybackslash}m{0.52\linewidth}>{\raggedleft\arraybackslash}m{0.27\linewidth}@{}}
\tikz\node[circle,fill=cor,text=white,minimum size=20pt,inner sep=0pt,font=\footnotesize\bfseries]{1}; &
{\color{cor}\large \faListUl} &
\textbf{Lista de funcionalidades (backlog)}\par{\scriptsize\color{cinza}\faUser\ Product Owner} &
\chip{\faCalendarCheck\ Semana 1}\par\vspace{2pt}\chip{\faClock\ cerca de 1h30}
\end{tabular}
\end{tcolorbox}
\begin{tcolorbox}[enhanced,colback=white,colframe=cor!25,boxrule=0.6pt,arc=6pt,left=6pt,right=8pt,top=4pt,bottom=4pt,before skip=0pt,after skip=4pt]
\begin{tabular}{@{}m{24pt}m{20pt}>{\raggedright\arraybackslash}m{0.52\linewidth}>{\raggedleft\arraybackslash}m{0.27\linewidth}@{}}
\tikz\node[circle,fill=cor,text=white,minimum size=20pt,inner sep=0pt,font=\footnotesize\bfseries]{2}; &
{\color{cor}\large \faColumns} &
\textbf{Quadro Kanban}\par{\scriptsize\color{cinza}\faUser\ Scrum Master} &
\chip{\faCalendarCheck\ Semana 1}\par\vspace{2pt}\chip{\faClock\ cerca de 45 min}
\end{tabular}
\end{tcolorbox}
\begin{tcolorbox}[enhanced,colback=suave!60,colframe=cor!25,boxrule=0.6pt,arc=6pt,left=6pt,right=8pt,top=4pt,bottom=4pt,before skip=0pt,after skip=4pt]
\begin{tabular}{@{}m{24pt}m{20pt}>{\raggedright\arraybackslash}m{0.52\linewidth}>{\raggedleft\arraybackslash}m{0.27\linewidth}@{}}
\tikz\node[circle,fill=cor,text=white,minimum size=20pt,inner sep=0pt,font=\footnotesize\bfseries]{3}; &
{\color{cor}\large \faRunning} &
\textbf{Planejamento de uma sprint}\par{\scriptsize\color{cinza}\faUser\ PO escolhe · Scrum Master calcula} &
\chip{\faCalendarCheck\ Semana 2}\par\vspace{2pt}\chip{\faClock\ cerca de 1h}
\end{tabular}
\end{tcolorbox}
\begin{tcolorbox}[enhanced,colback=white,colframe=cor!25,boxrule=0.6pt,arc=6pt,left=6pt,right=8pt,top=4pt,bottom=4pt,before skip=0pt,after skip=4pt]
\begin{tabular}{@{}m{24pt}m{20pt}>{\raggedright\arraybackslash}m{0.52\linewidth}>{\raggedleft\arraybackslash}m{0.27\linewidth}@{}}
\tikz\node[circle,fill=cor,text=white,minimum size=20pt,inner sep=0pt,font=\footnotesize\bfseries]{4}; &
{\color{cor}\large \faComments} &
\textbf{Reuniões rápidas e retrospectiva}\par{\scriptsize\color{cinza}\faUser\ Scrum Master} &
\chip{\faCalendarCheck\ Semanas 2 e 3}\par\vspace{2pt}\chip{\faClock\ 15 min por reunião}
\end{tabular}
\end{tcolorbox}
\secao{\faInfoCircle}{Como usar este documento}
\begin{enumerate}
  \item Sigam as etapas na ordem. Cada uma diz \textbf{em que semana} fazer, \textbf{quanto tempo} leva e qual \textbf{aula da apostila} explica o assunto.
  \item Nas tabelas, a \textbf{1ª linha é um exemplo}: usem como modelo para preencher as outras.
  \item O \textbf{Desafio extra} é opcional: só para quem quiser ir além.
  \item Preencham no computador pela versão \textbf{Word}; este PDF é para ler e imprimir (há espaço para escrever à mão).
  \item Só considerem a etapa concluída quando todos os quadradinhos do \textbf{Pronto quando} estiverem marcados.
  \item Ao terminar, enviem o arquivo pelo site da divisão de funções, em \textbf{Entregas da Frente B}, com o nome do selo \textbf{Entregar no site}.
\end{enumerate}
\clearpage
\noindent\begin{tikzpicture}
\draw[cor!25,line width=2.2pt](0.00,0)--(1.50,0);
\draw[cor!25,line width=2.2pt](1.50,0)--(3.00,0);
\draw[cor!25,line width=2.2pt](3.00,0)--(4.50,0);
\node[circle,fill=cor,text=white,draw=suave,line width=3pt,minimum size=24pt,inner sep=0pt,font=\small\bfseries] at (0.00,0){1};
\node[circle,fill=white,draw=cor!40,line width=1.2pt,text=cor!70,minimum size=16pt,inner sep=0pt,font=\scriptsize\bfseries] at (1.50,0){2};
\node[circle,fill=white,draw=cor!40,line width=1.2pt,text=cor!70,minimum size=16pt,inner sep=0pt,font=\scriptsize\bfseries] at (3.00,0){3};
\node[circle,fill=white,draw=cor!40,line width=1.2pt,text=cor!70,minimum size=16pt,inner sep=0pt,font=\scriptsize\bfseries] at (4.50,0){4};
\end{tikzpicture}
\par\vspace{8pt}
\begin{faixa}[before skip=0pt]
\begin{tabular}{@{}m{40pt}>{\raggedright\arraybackslash}m{0.78\linewidth}@{}}
\tikz\node[circle,fill=white,text=cor,minimum size=36pt,inner sep=0pt]{\Large \faListUl}; &
{\footnotesize\bfseries ETAPA 1 DE 4}\ \ \chipbranco{\faCalendarCheck\ Semana 1}\ \chipbranco{\faClock\ cerca de 1h30}\par\vspace{2pt}
{\fontsize{17}{21}\selectfont\bfseries Lista de funcionalidades (backlog)}\par\vspace{3pt}
{\footnotesize\faUser\ \ Quem faz: \textbf{Product Owner}\ \ ·\ \ \faCompass\ \ Ajuda: Guia da Frente B, manual da sua função}
\end{tabular}
\end{faixa}
\begin{apostila}\vspace{4pt}Aulas 01 e 02 · O Product Backlog é a lista de tudo que o produto precisa, em ordem de importância. Quem cuida dela é o Product Owner (PO). Cada item é escrito do ponto de vista de quem usa.\end{apostila}
\rotulo{\faBullseye}{Objetivo}
{\large Listar 10 coisas que o Estadia faz, escritas como histórias de usuário e em ordem de importância.}\par
\rotulo{\faShoePrints}{Passo a passo}
\begin{enumerate}
  \item Peçam ao dev de apoio para mostrar o sistema de demonstração (15 minutos bastam).
  \item Escrevam 10 itens no formato: ``Como \_\_\_, quero \_\_\_, para \_\_\_''. Veja o exemplo na tabela.
  \item Deem prioridade: Alta (sem isso o hotel não funciona), Média ou Baixa.
  \item Para cada item, escrevam como saber que está pronto (uma frase).
\end{enumerate}
\rotulo{\faPen}{Para preencher}
{\footnotesize\color{cinza}\faHandPointDown\ A 1ª linha é um exemplo. Preencham as outras.}\par\vspace{2pt}
\begingroup\small
\begin{longtable}{>{\raggedright\arraybackslash}p{0.042\linewidth}>{\raggedright\arraybackslash}p{0.329\linewidth}>{\raggedright\arraybackslash}p{0.156\linewidth}>{\raggedright\arraybackslash}p{0.329\linewidth}}
\rowcolor{cor} \color{white}\textbf{\footnotesize \#} & \color{white}\textbf{\footnotesize Como \_\_\_, quero \_\_\_, para \_\_\_} & \color{white}\textbf{\footnotesize Prioridade} & \color{white}\textbf{\footnotesize Está pronto quando…} \\
\endfirsthead
\rowcolor{cor} \color{white}\textbf{\footnotesize \#} & \color{white}\textbf{\footnotesize Como \_\_\_, quero \_\_\_, para \_\_\_} & \color{white}\textbf{\footnotesize Prioridade} & \color{white}\textbf{\footnotesize Está pronto quando…} \\
\endhead
\rowcolor{white}\textbf{1} & Como recepcionista, quero ver os quartos livres, para não vender um quarto ocupado & Alta & O sistema mostra só os quartos livres na data escolhida \\ \hline
\rowcolor{suave!55}\rule{0pt}{24pt}\textbf{2} &  &  &  \\ \hline
\rowcolor{white}\rule{0pt}{24pt}\textbf{3} &  &  &  \\ \hline
\rowcolor{suave!55}\rule{0pt}{24pt}\textbf{4} &  &  &  \\ \hline
\rowcolor{white}\rule{0pt}{24pt}\textbf{5} &  &  &  \\ \hline
\rowcolor{suave!55}\rule{0pt}{24pt}\textbf{6} &  &  &  \\ \hline
\rowcolor{white}\rule{0pt}{24pt}\textbf{7} &  &  &  \\ \hline
\rowcolor{suave!55}\rule{0pt}{24pt}\textbf{8} &  &  &  \\ \hline
\rowcolor{white}\rule{0pt}{24pt}\textbf{9} &  &  &  \\ \hline
\rowcolor{suave!55}\rule{0pt}{24pt}\textbf{10} &  &  &  \\ \hline
\end{longtable}
\endgroup
\begin{lrbox}{\blocofinal}\begin{minipage}{\linewidth}
\begin{tcolorbox}[enhanced,colback=suave!45,colframe=cor!35,boxrule=0.8pt,arc=8pt,left=10pt,right=10pt,top=8pt,bottom=8pt,before skip=0pt,after skip=0pt]
\begin{minipage}[t]{0.56\linewidth}\vspace{0pt}
{\footnotesize\bfseries\color{cor}\faCheckDouble\ \ PRONTO QUANDO}\par\vspace{2pt}
\begin{checklist}
  \item 10 itens escritos no formato certo.
  \item Todos com prioridade e ``está pronto quando''.
\end{checklist}
\end{minipage}\hfill
\begin{minipage}[t]{0.40\linewidth}\vspace{0pt}
\begin{selo}\centering
{\LARGE\color{cor}\faUpload}\par\vspace{3pt}
{\scriptsize\bfseries\color{cor}ENTREGAR NO SITE COMO}\par\vspace{2pt}
{\bfseries B1-backlog}
\end{selo}
\end{minipage}
\end{tcolorbox}
\end{minipage}\end{lrbox}
\usarblocofinal
\Needspace{6\baselineskip}
\begin{tcolorbox}[enhanced,colback=white,colframe=cor,boxrule=1pt,arc=8pt,left=12pt,right=12pt,top=12pt,bottom=8pt,before skip=20pt,
  overlay={\node[fill=cor,text=white,rounded corners=3pt,font=\scriptsize\bfseries,inner sep=4pt,anchor=west] at ([xshift=10pt]frame.north west) {\faMedal\ DESAFIO EXTRA · OPCIONAL};}]
Escrevam os critérios no formato da aula: ``Dado \_\_\_, quando \_\_\_, então \_\_\_''.
\end{tcolorbox}
\ifjaanotou\else\anotacoes\fi
\clearpage
\noindent\begin{tikzpicture}
\draw[cor,line width=2.2pt](0.00,0)--(1.50,0);
\draw[cor!25,line width=2.2pt](1.50,0)--(3.00,0);
\draw[cor!25,line width=2.2pt](3.00,0)--(4.50,0);
\node[circle,fill=cor,text=white,minimum size=16pt,inner sep=0pt,font=\scriptsize] at (0.00,0){\faCheck};
\node[circle,fill=cor,text=white,draw=suave,line width=3pt,minimum size=24pt,inner sep=0pt,font=\small\bfseries] at (1.50,0){2};
\node[circle,fill=white,draw=cor!40,line width=1.2pt,text=cor!70,minimum size=16pt,inner sep=0pt,font=\scriptsize\bfseries] at (3.00,0){3};
\node[circle,fill=white,draw=cor!40,line width=1.2pt,text=cor!70,minimum size=16pt,inner sep=0pt,font=\scriptsize\bfseries] at (4.50,0){4};
\end{tikzpicture}
\par\vspace{8pt}
\begin{faixa}[before skip=0pt]
\begin{tabular}{@{}m{40pt}>{\raggedright\arraybackslash}m{0.78\linewidth}@{}}
\tikz\node[circle,fill=white,text=cor,minimum size=36pt,inner sep=0pt]{\Large \faColumns}; &
{\footnotesize\bfseries ETAPA 2 DE 4}\ \ \chipbranco{\faCalendarCheck\ Semana 1}\ \chipbranco{\faClock\ cerca de 45 min}\par\vspace{2pt}
{\fontsize{17}{21}\selectfont\bfseries Quadro Kanban}\par\vspace{3pt}
{\footnotesize\faUser\ \ Quem faz: \textbf{Scrum Master}\ \ ·\ \ \faCompass\ \ Ajuda: Guia da Frente B, manual da sua função}
\end{tabular}
\end{faixa}
\begin{apostila}\vspace{4pt}Aula 04 · Kanban é um quadro com colunas (A fazer → Fazendo → Feito) e cartões que andam pelas colunas. Ajuda todo mundo a ver o que está sendo feito e evita começar coisas demais ao mesmo tempo.\end{apostila}
\rotulo{\faBullseye}{Objetivo}
{\large Criar um quadro onde o grupo inteiro vê o andamento das tarefas.}\par
\rotulo{\faShoePrints}{Passo a passo}
\begin{enumerate}
  \item Criem um quadro gratuito no Trello com 3 colunas: A fazer, Fazendo, Feito.
  \item Criem um cartão para cada entrega das 4 frentes (estão no site), com uma etiqueta de cor por frente.
  \item Convidem o grupo e combinem: cada um move o próprio cartão.
  \item Combinem um limite: no máximo 2 cartões em ``Fazendo'' por pessoa.
  \item Tirem um print agora e outro na última semana.
\end{enumerate}
\rotulo{\faPen}{Para preencher}
\Needspace{14\baselineskip}
\begin{tcolorbox}[enhanced,colback=white,colframe=cor!30,boxrule=0.8pt,arc=5pt,left=8pt,right=8pt,top=4pt,bottom=6pt,before skip=0pt,after skip=5pt]
{\scriptsize\bfseries\color{cor}\MakeUppercase{Link do quadro}}\par\vspace{16pt}\noindent\tikz\draw[cor!45,dotted,line width=1.3pt](0,0)--(\linewidth,0);\par
\end{tcolorbox}
\begin{tcolorbox}[enhanced,colback=white,colframe=cor!30,boxrule=0.8pt,arc=5pt,left=8pt,right=8pt,top=4pt,bottom=6pt,before skip=0pt,after skip=5pt]
{\scriptsize\bfseries\color{cor}\MakeUppercase{Cores usadas para cada frente}}\par\vspace{16pt}\noindent\tikz\draw[cor!45,dotted,line width=1.3pt](0,0)--(\linewidth,0);\par
\end{tcolorbox}
\begin{tcolorbox}[enhanced,colback=white,colframe=cor!30,boxrule=0.8pt,arc=5pt,left=8pt,right=8pt,top=4pt,bottom=6pt,before skip=0pt,after skip=5pt]
{\scriptsize\bfseries\color{cor}\MakeUppercase{Limite de cartões em ``Fazendo''}}\par\vspace{16pt}\noindent\tikz\draw[cor!45,dotted,line width=1.3pt](0,0)--(\linewidth,0);\par
\end{tcolorbox}
\vspace{4pt}
\begin{lrbox}{\blocofinal}\begin{minipage}{\linewidth}
\begin{tcolorbox}[enhanced,colback=suave!45,colframe=cor!35,boxrule=0.8pt,arc=8pt,left=10pt,right=10pt,top=8pt,bottom=8pt,before skip=0pt,after skip=0pt]
\begin{minipage}[t]{0.56\linewidth}\vspace{0pt}
{\footnotesize\bfseries\color{cor}\faCheckDouble\ \ PRONTO QUANDO}\par\vspace{2pt}
\begin{checklist}
  \item Quadro com as entregas das 4 frentes.
  \item Prints do começo e do fim.
\end{checklist}
\end{minipage}\hfill
\begin{minipage}[t]{0.40\linewidth}\vspace{0pt}
\begin{selo}\centering
{\LARGE\color{cor}\faUpload}\par\vspace{3pt}
{\scriptsize\bfseries\color{cor}ENTREGAR NO SITE COMO}\par\vspace{2pt}
{\bfseries B2-kanban (prints)}
\end{selo}
\end{minipage}
\end{tcolorbox}
\end{minipage}\end{lrbox}
\usarblocofinal
\ifjaanotou\else\anotacoes\fi
\clearpage
\noindent\begin{tikzpicture}
\draw[cor,line width=2.2pt](0.00,0)--(1.50,0);
\draw[cor,line width=2.2pt](1.50,0)--(3.00,0);
\draw[cor!25,line width=2.2pt](3.00,0)--(4.50,0);
\node[circle,fill=cor,text=white,minimum size=16pt,inner sep=0pt,font=\scriptsize] at (0.00,0){\faCheck};
\node[circle,fill=cor,text=white,minimum size=16pt,inner sep=0pt,font=\scriptsize] at (1.50,0){\faCheck};
\node[circle,fill=cor,text=white,draw=suave,line width=3pt,minimum size=24pt,inner sep=0pt,font=\small\bfseries] at (3.00,0){3};
\node[circle,fill=white,draw=cor!40,line width=1.2pt,text=cor!70,minimum size=16pt,inner sep=0pt,font=\scriptsize\bfseries] at (4.50,0){4};
\end{tikzpicture}
\par\vspace{8pt}
\begin{faixa}[before skip=0pt]
\begin{tabular}{@{}m{40pt}>{\raggedright\arraybackslash}m{0.78\linewidth}@{}}
\tikz\node[circle,fill=white,text=cor,minimum size=36pt,inner sep=0pt]{\Large \faRunning}; &
{\footnotesize\bfseries ETAPA 3 DE 4}\ \ \chipbranco{\faCalendarCheck\ Semana 2}\ \chipbranco{\faClock\ cerca de 1h}\par\vspace{2pt}
{\fontsize{17}{21}\selectfont\bfseries Planejamento de uma sprint}\par\vspace{3pt}
{\footnotesize\faUser\ \ Quem faz: \textbf{PO escolhe · Scrum Master calcula}\ \ ·\ \ \faCompass\ \ Ajuda: Guia da Frente B, manual da sua função}
\end{tabular}
\end{faixa}
\begin{apostila}\vspace{4pt}Aulas 01 e 02 · Sprint é um ciclo curto de trabalho com um objetivo claro. No começo, o time escolhe o que cabe no tempo que tem. Capacidade = pessoas × horas por dia × dias.\end{apostila}
\rotulo{\faBullseye}{Objetivo}
{\large Planejar uma semana de trabalho do grupo do jeito Scrum.}\par
\rotulo{\faShoePrints}{Passo a passo}
\begin{enumerate}
  \item Escrevam o objetivo da sprint em uma frase. Ex.: ``Terminar a primeira versão de todas as entregas''.
  \item Calculem a capacidade: quantas pessoas × horas por dia × dias. Ex.: 8 pessoas × 1 h × 5 dias = 40 horas.
  \item Escolham do quadro Kanban as tarefas que cabem nessas horas e estimem cada uma.
  \item Se passar da capacidade, tirem tarefas da sprint.
\end{enumerate}
\rotulo{\faPen}{Para preencher}
{\footnotesize\color{cinza}\faHandPointDown\ A 1ª linha é um exemplo. Preencham as outras.}\par\vspace{2pt}
\begingroup\small
\begin{longtable}{>{\raggedright\arraybackslash}p{0.297\linewidth}>{\raggedright\arraybackslash}p{0.297\linewidth}>{\raggedright\arraybackslash}p{0.297\linewidth}}
\rowcolor{cor} \color{white}\textbf{\footnotesize Tarefa} & \color{white}\textbf{\footnotesize Quem faz} & \color{white}\textbf{\footnotesize Horas estimadas} \\
\endfirsthead
\rowcolor{cor} \color{white}\textbf{\footnotesize Tarefa} & \color{white}\textbf{\footnotesize Quem faz} & \color{white}\textbf{\footnotesize Horas estimadas} \\
\endhead
\rowcolor{white}\textbf{Canvas do projeto} & Frente A & 3 \\ \hline
\rowcolor{suave!55}\rule{0pt}{24pt} &  &  \\ \hline
\rowcolor{white}\rule{0pt}{24pt} &  &  \\ \hline
\rowcolor{suave!55}\rule{0pt}{24pt} &  &  \\ \hline
\rowcolor{white}\rule{0pt}{24pt} &  &  \\ \hline
\rowcolor{suave!55}\rule{0pt}{24pt} &  &  \\ \hline
\end{longtable}
\endgroup
\subsecao{\faEdit\ Resumo da sprint}
\Needspace{14\baselineskip}
\begin{tcolorbox}[enhanced,colback=white,colframe=cor!30,boxrule=0.8pt,arc=5pt,left=8pt,right=8pt,top=4pt,bottom=6pt,before skip=0pt,after skip=5pt]
{\scriptsize\bfseries\color{cor}\MakeUppercase{Objetivo da sprint}}\par\vspace{16pt}\noindent\tikz\draw[cor!45,dotted,line width=1.3pt](0,0)--(\linewidth,0);\par
\end{tcolorbox}
\begin{tcolorbox}[enhanced,colback=white,colframe=cor!30,boxrule=0.8pt,arc=5pt,left=8pt,right=8pt,top=4pt,bottom=6pt,before skip=0pt,after skip=5pt]
{\scriptsize\bfseries\color{cor}\MakeUppercase{Capacidade (conta)}}\par\vspace{16pt}\noindent\tikz\draw[cor!45,dotted,line width=1.3pt](0,0)--(\linewidth,0);\par
\end{tcolorbox}
\begin{tcolorbox}[enhanced,colback=white,colframe=cor!30,boxrule=0.8pt,arc=5pt,left=8pt,right=8pt,top=4pt,bottom=6pt,before skip=0pt,after skip=5pt]
{\scriptsize\bfseries\color{cor}\MakeUppercase{Total de horas planejadas}}\par\vspace{16pt}\noindent\tikz\draw[cor!45,dotted,line width=1.3pt](0,0)--(\linewidth,0);\par
\end{tcolorbox}
\vspace{4pt}
\begin{lrbox}{\blocofinal}\begin{minipage}{\linewidth}
\begin{tcolorbox}[enhanced,colback=suave!45,colframe=cor!35,boxrule=0.8pt,arc=8pt,left=10pt,right=10pt,top=8pt,bottom=8pt,before skip=0pt,after skip=0pt]
\begin{minipage}[t]{0.56\linewidth}\vspace{0pt}
{\footnotesize\bfseries\color{cor}\faCheckDouble\ \ PRONTO QUANDO}\par\vspace{2pt}
\begin{checklist}
  \item Objetivo em uma frase.
  \item Total planejado menor ou igual à capacidade.
\end{checklist}
\end{minipage}\hfill
\begin{minipage}[t]{0.40\linewidth}\vspace{0pt}
\begin{selo}\centering
{\LARGE\color{cor}\faUpload}\par\vspace{3pt}
{\scriptsize\bfseries\color{cor}ENTREGAR NO SITE COMO}\par\vspace{2pt}
{\bfseries B3-sprint}
\end{selo}
\end{minipage}
\end{tcolorbox}
\end{minipage}\end{lrbox}
\usarblocofinal
\Needspace{6\baselineskip}
\begin{tcolorbox}[enhanced,colback=white,colframe=cor,boxrule=1pt,arc=8pt,left=12pt,right=12pt,top=12pt,bottom=8pt,before skip=20pt,
  overlay={\node[fill=cor,text=white,rounded corners=3pt,font=\scriptsize\bfseries,inner sep=4pt,anchor=west] at ([xshift=10pt]frame.north west) {\faMedal\ DESAFIO EXTRA · OPCIONAL};}]
Escrevam a ``Definição de Pronto'' (DoD) do grupo: o que toda entrega precisa ter para ser considerada pronta.
\end{tcolorbox}
\ifjaanotou\else\anotacoes\fi
\clearpage
\noindent\begin{tikzpicture}
\draw[cor,line width=2.2pt](0.00,0)--(1.50,0);
\draw[cor,line width=2.2pt](1.50,0)--(3.00,0);
\draw[cor,line width=2.2pt](3.00,0)--(4.50,0);
\node[circle,fill=cor,text=white,minimum size=16pt,inner sep=0pt,font=\scriptsize] at (0.00,0){\faCheck};
\node[circle,fill=cor,text=white,minimum size=16pt,inner sep=0pt,font=\scriptsize] at (1.50,0){\faCheck};
\node[circle,fill=cor,text=white,minimum size=16pt,inner sep=0pt,font=\scriptsize] at (3.00,0){\faCheck};
\node[circle,fill=cor,text=white,draw=suave,line width=3pt,minimum size=24pt,inner sep=0pt,font=\small\bfseries] at (4.50,0){4};
\end{tikzpicture}
\par\vspace{8pt}
\begin{faixa}[before skip=0pt]
\begin{tabular}{@{}m{40pt}>{\raggedright\arraybackslash}m{0.78\linewidth}@{}}
\tikz\node[circle,fill=white,text=cor,minimum size=36pt,inner sep=0pt]{\Large \faComments}; &
{\footnotesize\bfseries ETAPA 4 DE 4}\ \ \chipbranco{\faCalendarCheck\ Semanas 2 e 3}\ \chipbranco{\faClock\ 15 min por reunião}\par\vspace{2pt}
{\fontsize{17}{21}\selectfont\bfseries Reuniões rápidas e retrospectiva}\par\vspace{3pt}
{\footnotesize\faUser\ \ Quem faz: \textbf{Scrum Master}\ \ ·\ \ \faCompass\ \ Ajuda: Guia da Frente B, manual da sua função}
\end{tabular}
\end{faixa}
\begin{apostila}\vspace{4pt}Aulas 01 e 02 · A daily é uma reunião de 15 minutos com 3 perguntas: o que fiz, o que vou fazer, tem algo me atrapalhando? A retrospectiva, no fim da sprint, pergunta o que melhorar.\end{apostila}
\rotulo{\faBullseye}{Objetivo}
{\large Conduzir as reuniões rápidas do grupo e registrar as melhorias.}\par
\rotulo{\faShoePrints}{Passo a passo}
\begin{enumerate}
  \item O Scrum Master puxa 2 dailies de 15 minutos (pode ser por mensagem no grupo).
  \item Anotem as respostas e qualquer impedimento.
  \item No fim da semana, façam a retrospectiva: o que foi bem, o que melhorar e 3 ações.
\end{enumerate}
\rotulo{\faPen}{Para preencher}
{\footnotesize\color{cinza}\faHandPointDown\ A 1ª linha é um exemplo. Preencham as outras.}\par\vspace{2pt}
\begingroup\small
\begin{longtable}{>{\raggedright\arraybackslash}p{0.069\linewidth}>{\raggedright\arraybackslash}p{0.256\linewidth}>{\raggedright\arraybackslash}p{0.256\linewidth}>{\raggedright\arraybackslash}p{0.275\linewidth}}
\rowcolor{cor} \color{white}\textbf{\footnotesize Reunião} & \color{white}\textbf{\footnotesize Data} & \color{white}\textbf{\footnotesize Principais pontos} & \color{white}\textbf{\footnotesize Impedimentos} \\
\endfirsthead
\rowcolor{cor} \color{white}\textbf{\footnotesize Reunião} & \color{white}\textbf{\footnotesize Data} & \color{white}\textbf{\footnotesize Principais pontos} & \color{white}\textbf{\footnotesize Impedimentos} \\
\endhead
\rowcolor{white}\textbf{Daily} & 28/10 & Frente C terminou os OKRs & Ninguém tem acesso ao Trello \\ \hline
\rowcolor{suave!55}\rule{0pt}{24pt}\textbf{Daily} &  &  &  \\ \hline
\rowcolor{white}\rule{0pt}{24pt}\textbf{Daily} &  &  &  \\ \hline
\end{longtable}
\endgroup
\subsecao{\faEdit\ Retrospectiva}
\begingroup\small
\begin{longtable}{>{\raggedright\arraybackslash}p{0.297\linewidth}>{\raggedright\arraybackslash}p{0.297\linewidth}>{\raggedright\arraybackslash}p{0.297\linewidth}}
\rowcolor{cor} \color{white}\textbf{\footnotesize O que foi bem} & \color{white}\textbf{\footnotesize O que melhorar} & \color{white}\textbf{\footnotesize Ação (com responsável)} \\
\endfirsthead
\rowcolor{cor} \color{white}\textbf{\footnotesize O que foi bem} & \color{white}\textbf{\footnotesize O que melhorar} & \color{white}\textbf{\footnotesize Ação (com responsável)} \\
\endhead
\rowcolor{white}\rule{0pt}{24pt} &  &  \\ \hline
\rowcolor{suave!55}\rule{0pt}{24pt} &  &  \\ \hline
\rowcolor{white}\rule{0pt}{24pt} &  &  \\ \hline
\end{longtable}
\endgroup
\begin{lrbox}{\blocofinal}\begin{minipage}{\linewidth}
\begin{tcolorbox}[enhanced,colback=suave!45,colframe=cor!35,boxrule=0.8pt,arc=8pt,left=10pt,right=10pt,top=8pt,bottom=8pt,before skip=0pt,after skip=0pt]
\begin{minipage}[t]{0.56\linewidth}\vspace{0pt}
{\footnotesize\bfseries\color{cor}\faCheckDouble\ \ PRONTO QUANDO}\par\vspace{2pt}
\begin{checklist}
  \item 2 dailies registradas.
  \item Retrospectiva com 3 ações.
\end{checklist}
\end{minipage}\hfill
\begin{minipage}[t]{0.40\linewidth}\vspace{0pt}
\begin{selo}\centering
{\LARGE\color{cor}\faUpload}\par\vspace{3pt}
{\scriptsize\bfseries\color{cor}ENTREGAR NO SITE COMO}\par\vspace{2pt}
{\bfseries B4-reunioes}
\end{selo}
\end{minipage}
\end{tcolorbox}
\end{minipage}\end{lrbox}
\usarblocofinal
\ifjaanotou\else\anotacoes\fi
\clearpage
\secao{\faLaptopCode}{O papel do dev de apoio}
\begin{itemize}
  \item Mostrar o sistema funcionando e criar usuários de teste.
  \item Conferir o que a frente escreveu sobre o sistema.
  \item Passar números e prints da versão de demonstração.
  \item Explicar o que for técnico em palavras simples.
\end{itemize}
\secao{\faUpload}{Como entregar no site}
\begin{enumerate}
  \item Abram o link da divisão de funções que o grupo usa.
  \item Na Frente B, toquem em \textbf{Enviar entrega} (ou arrastem o arquivo): Word, PDF, Excel, PowerPoint, imagem ou ZIP, até 25 MB.
  \item O arquivo aparece na hora para todo o grupo, com o nome de quem enviou e a data. Versões novas podem ser enviadas: a mais recente fica no topo.
  \item Sem o servidor ligado (página aberta como arquivo), o envio não funciona: mandem o arquivo para a Frente A.
\end{enumerate}
\vspace{16pt}
\begin{faixa}
\centering{\Large\faTrophy}\par\vspace{4pt}{\large\bfseries Terminaram tudo? Avisem a Frente A e comemorem!}
\end{faixa}
\end{document}
